% EOS Methods Writeup — AASTeX63 format
% Drop this file into an Overleaf AASTeX project
\documentclass[twocolumn]{aastex701}

\usepackage{amsmath,amssymb}
\usepackage{bm}

\shorttitle{EOS Mixing and Inversion Methods}
\shortauthors{Tejada Arevalo}

\begin{document}

\title{Equation of State Framework: Volume Addition Law, Inversions, and Thermodynamic Derivatives}

\author{Roberto Tejada Arevalo}
\email{rarevalo@ucsc.edu}
\affiliation{University of California, Santa Cruz}

\begin{abstract}
We describe the computational framework used to construct multi-component equations of state (EOS) for planetary interior modeling.
The framework combines hydrogen--helium EOS tables with heavy-element EOS data via the volume addition law (VAL), incorporates ideal and non-ideal entropy of mixing corrections, performs numerical inversions between thermodynamic coordinate systems using hybrid root-finding algorithms, and computes thermodynamic derivatives relevant to convective stability.
\end{abstract}

\keywords{planets and satellites: interiors --- equation of state --- methods: numerical}

%============================================================
\section{Introduction}\label{sec:intro}
%============================================================

Modeling planetary interiors requires an equation of state that spans a wide range of pressures ($10^{6}$--$10^{14}\;\mathrm{dyn\;cm^{-2}}$), temperatures ($10^{2}$--$10^{5}\;\mathrm{K}$), and compositions.
We parameterize the composition by the helium mass fraction $Y$ and the total heavy-element (metal) mass fraction $Z$, so that the hydrogen mass fraction is $X = 1 - Y - Z$.
The heavy-element component may itself be a mixture of water, methane, ammonia, silicates (post-perovskite, MgSiO$_3$), and iron.

We adopt pre-computed EOS tables for pure hydrogen and pure helium \citep[e.g.,][]{Chabrier2019,Chabrier2021} and for heavy elements such as the AQUA water EOS \citep{Haldemann2020}, and combine them using the volume addition law.
Non-ideal mixing corrections for the H--He system follow \citet{Howard2023a}.

In the following sections we describe (1) the volume addition law and entropy of mixing (\S\ref{sec:val}), (2) the interpolation architecture (\S\ref{sec:interp}), (3) the inversion procedure for changing thermodynamic basis (\S\ref{sec:inversion}), (4) the derived thermodynamic quantities (\S\ref{sec:derivatives}), and (5) the ideal gas EOS with molecular rotational and vibrational degrees of freedom (\S\ref{sec:ideal_eos}).

%============================================================
\section{Volume Addition Law}\label{sec:val}
%============================================================

The volume addition law (VAL) constructs a multi-component EOS by assuming that each species occupies its pure-substance specific volume at the mixture pressure and temperature.
We work with mass-specific (per gram) quantities throughout.

\subsection{Notation}\label{sec:notation}

We define the \emph{helium prime fraction} $Y' \equiv Y/(1-Z)$, which is the helium mass fraction within the hydrogen--helium sub-mixture.
Quantities subscripted ``$xy$'' refer to the H--He component evaluated at $(P,T,Y')$; quantities subscripted ``$z$'' refer to the heavy-element component evaluated at $(P,T)$.

\subsection{Density}\label{sec:val_rho}

The specific volume of the mixture is
\begin{equation}\label{eq:val_rho}
    v_{\rm mix} = v_{xy}(P,T,Y')\,(1-Z) + \frac{Z}{\rho_z(P,T)},
\end{equation}
where the H--He specific volume is itself constructed from the pure-species volumes:
\begin{equation}\label{eq:vxy}
    v_{xy} = \frac{1-Y'}{\rho_{\rm H}(P,T)} + \frac{Y'}{\rho_{\rm He}(P,T)} + \Delta v_{\rm mix}(P,T,Y').
\end{equation}
Here $\Delta v_{\rm mix}$ is the non-ideal volume of mixing from \citet{Howard2023a}, parameterized as
\begin{equation}\label{eq:vmix}
    \Delta v_{\rm mix}(P,T,Y') = \widetilde{V}_{\rm mix}(P,T)\,(1-Y')\,Y',
\end{equation}
where $\widetilde{V}_{\rm mix}(P,T)$ is interpolated from tabulated interaction data.
When the \citet{Chabrier2021} (CD21) tables are used, the H--He interaction terms are already folded into the tables and $\Delta v_{\rm mix} = 0$.
The mixture density follows as $\rho_{\rm mix} = 1/v_{\rm mix}$.

For heavy-element mixtures (e.g., ice--rock), the heavy-element density is itself obtained by volume addition:
\begin{equation}
    \frac{1}{\rho_z} = \frac{1-f_{\rm rock}}{\rho_{\rm ice}(P,T)} + \frac{f_{\rm rock}}{\rho_{\rm rock}(P,T)},
\end{equation}
where $f_{\rm rock}$ is the rock mass fraction within the heavy-element component.

\subsection{Entropy}\label{sec:val_s}

The mixture entropy per unit mass is
\begin{equation}\label{eq:val_s}
    s_{\rm mix} = s_{xy}\,(1-Z) + s_z\,Z + s_{\rm mix,id} + \Delta s_{\rm mix,nonid}\,(1-Z),
\end{equation}
where $s_{xy}$ is the H--He entropy constructed from pure-species entropies (without the ideal mixing term):
\begin{equation}\label{eq:sxy}
    s_{xy} = s_{\rm H}(P,T)\,(1-Y') + s_{\rm He}(P,T)\,Y'.
\end{equation}
The ideal entropy of mixing $s_{\rm mix,id}$ and the non-ideal correction $\Delta s_{\rm mix,nonid}$ are described in \S\ref{sec:smix}.

\subsection{Internal Energy}\label{sec:val_u}

The mixture internal energy per unit mass is
\begin{equation}\label{eq:val_u}
    u_{\rm mix} = u_{xy}(P,T,Y')\,(1-Z) + u_z(P,T)\,Z,
\end{equation}
where
\begin{equation}
    u_{xy} = u_{\rm H}(P,T)\,(1-Y') + u_{\rm He}(P,T)\,Y' + \Delta u_{\rm mix}(P,T,Y'),
\end{equation}
and $\Delta u_{\rm mix} = \widetilde{U}_{\rm mix}(P,T)\,(1-Y')\,Y'$ is the non-ideal internal energy of mixing \citep{Howard2023a}.

\subsection{Entropy of Mixing}\label{sec:smix}

\subsubsection{Ideal Entropy of Mixing}

For a mixture of $N$ molecular species with mass fractions $\{f_i\}$ and molecular weights $\{\mu_i\}$, the ideal entropy of mixing per unit mass (in units of $k_B$ per baryon) is
\begin{equation}\label{eq:smix_id}
    s_{\rm mix,id} = -\frac{1}{\bar{\mu}} \sum_{i=1}^{N} x_i \ln x_i,
\end{equation}
where the number fractions are
\begin{equation}
    x_i = \frac{f_i / \mu_i}{\sum_j f_j / \mu_j},
\end{equation}
and the mean molecular weight is
\begin{equation}
    \bar{\mu} = \sum_{i=1}^{N} \mu_i\, x_i.
\end{equation}
Our framework supports up to seven species: H ($\mu=1$), He ($\mu=4.003$), H$_2$O ($\mu=18.015$), CH$_4$ ($\mu=16.04$), NH$_3$ ($\mu=17.031$), MgSiO$_3$ ($\mu=100.389$), and Fe ($\mu=55.845$).
When a species has zero mass fraction, we use a guarded logarithm $x_i \ln x_i \to 0$ in the limit $x_i \to 0$.

For the special case of a binary H--He mixture, this reduces to
\begin{equation}\label{eq:smix_id_y}
    s_{\rm mix,id}^{(xy)} = -\frac{x_{\rm H}\ln x_{\rm H} + x_{\rm He}\ln x_{\rm He}}{\mu_{\rm H}\,x_{\rm H} + \mu_{\rm He}\,x_{\rm He}},
\end{equation}
where $x_{\rm He} = Y'/\mu_{\rm He}\,/\,(Y'/\mu_{\rm He} + (1-Y')/\mu_{\rm H})$.

\subsubsection{Non-Ideal Entropy of Mixing}

The non-ideal correction to the H--He entropy of mixing is
\begin{equation}\label{eq:smix_nonid}
    \Delta s_{\rm mix,nonid} = \widetilde{S}_{\rm mix}(P,T)\,(1-Y')\,Y' - s_{\rm mix,id}^{(xy)},
\end{equation}
where $\widetilde{S}_{\rm mix}(P,T)$ is interpolated from the tabulated total (ideal + non-ideal) entropy of mixing of \citet{Howard2023a}.
This correction is applied only when the CMS19 H--He tables are used; for the CD21 tables the interaction terms are already included.

\subsubsection{Composition Parameterization}

The heavy-element mass fractions may be specified in two ways:
\begin{enumerate}
    \item \textbf{Nested conditional fractions:} The mass fractions are defined hierarchically as $f_{\rm H} = (1-Y')(1-Z_w)(1-Z_m)(1-Z_a)(1-Z_r)(1-Z_{\rm Fe})$, $f_{\rm He} = Y'(1-Z_w)\cdots$, $f_{\rm H_2O} = Z_w(1-Z_m)\cdots$, etc.
    \item \textbf{Absolute (physical) mass fractions:} The mass fractions $\{f_i\}$ are specified directly and must satisfy $\sum_i f_i = 1$.
\end{enumerate}
The first parameterization is natural for the nested VAL construction; the second is used for ice--rock mixtures where the physical fractions within the $Z$ component are specified independently.

%============================================================
\section{Interpolation Architecture}\label{sec:interp}
%============================================================

\subsection{Table Structure}

Thermodynamic quantities are stored on regular grids and accessed via multi-linear interpolation using \texttt{scipy.interpolate.RegularGridInterpolator}.
The tables are four-dimensional, with coordinates depending on the thermodynamic basis:

\begin{enumerate}
    \item \textbf{$P$--$T$ tables:} Coordinates $(\log P,\;\log T,\;Y,\;Z)$.
    Dependent variables: $s(P,T,Y,Z)$, $\log\rho(P,T,Y,Z)$, $\log u(P,T,Y,Z)$.
    Units: $s$ in $\mathrm{erg\;g^{-1}\;K^{-1}}$, $\rho$ in $\mathrm{g\;cm^{-3}}$, $u$ in $\mathrm{erg\;g^{-1}}$, $P$ in $\mathrm{dyn\;cm^{-2}}$.

    \item \textbf{$\rho$--$T$ tables:} Coordinates $(\log\rho,\;\log T,\;Y,\;Z)$.
    Dependent variables: $s(\rho,T,Y,Z)$, $\log P(\rho,T,Y,Z)$.

    \item \textbf{$S$--$P$ tables:} Coordinates $(s,\;\log P,\;Y,\;Z)$, with $s$ in $k_B\;\mathrm{baryon}^{-1}$.
    Dependent variables: $\log T(s,P,Y,Z)$, $\log\rho(s,P,Y,Z)$.

    \item \textbf{$S$--$\rho$ tables:} Coordinates $(s,\;\log\rho,\;Y,\;Z)$.
    Dependent variables: $\log P(s,\rho,Y,Z)$, $\log T(s,\rho,Y,Z)$.
\end{enumerate}

The $P$--$T$ tables are the primary (native) representation and are constructed directly from the VAL.
The remaining tables are produced by numerical inversion of the $P$--$T$ tables (\S\ref{sec:inversion}).

\subsection{Pure-Species Interpolators}

For the CMS19 and CD21 H--He EOS, the pure hydrogen and pure helium tables are stored separately as 2D grids on $(\log P, \log T)$ and interpolated independently before being combined via the VAL.
Non-ideal mixing quantities $\widetilde{S}_{\rm mix}$, $\widetilde{V}_{\rm mix}$, and $\widetilde{U}_{\rm mix}$ from \citet{Howard2023a} are likewise stored on 2D $(\log P, \log T)$ grids.

\subsection{Multi-Fraction Extension}

For models requiring a variable rock-to-ice ratio, the tables are extended to five dimensions by stacking pre-computed 4D tables at discrete rock fractions $f_{\rm ppv} \in \{0.0, 0.25, 0.5, 0.75, 1.0\}$.
The resulting interpolators have coordinates $(\log P, \log T, Y, Z, f_{\rm ppv})$ and support continuous interpolation in all five dimensions.

%============================================================
\section{EOS Inversions}\label{sec:inversion}
%============================================================

\subsection{Motivation}

Stellar and planetary structure codes require the EOS in different thermodynamic bases depending on the formulation of the structure equations.
For example, adiabatic interior models naturally work in entropy--pressure $(s,P)$ coordinates, while hydrodynamic simulations may require density--temperature $(\rho,T)$ or entropy--density $(s,\rho)$ representations.

Since the VAL provides thermodynamic quantities as functions of $(P,T)$, obtaining quantities in other coordinate systems requires numerical inversion.
We solve for the dependent variables by root-finding.

\subsection{Inversion Strategy}\label{sec:inv_strategy}

Consider the generic problem of finding $\log T$ given $(s, \log P, Y, Z)$.
We define the error function
\begin{equation}\label{eq:err_sp}
    \varepsilon(\log T) = \frac{s_{\rm VAL}(\log P, \log T, Y, Z)}{s_{\rm target}} - 1,
\end{equation}
and seek $\log T^*$ such that $\varepsilon(\log T^*) = 0$.

The initial guess is obtained from an ideal-gas EOS or, when available, from the solution at the previous composition point to accelerate convergence across the $Z$ grid.

\subsubsection{Hybrid Newton--Brentq Method}

Our default solver (``newton\_brentq'') proceeds as follows:
\begin{enumerate}
    \item Attempt Newton's method with tolerance $10^{-5}$ and a maximum of 100 iterations.
    \item If Newton's method fails to converge (raises a \texttt{RuntimeError}), fall back to Brent's method:
    \begin{enumerate}
        \item Initialize a bracketing interval $[a,b] = [\log T_{\rm guess} - \delta,\;\log T_{\rm guess} + \delta]$ with $\delta = 0.1$.
        \item If $\varepsilon(a)\,\varepsilon(b) < 0$, apply \texttt{brentq} with tolerance $10^{-5}$.
        \item Otherwise, expand the interval by a factor of 2 and retry, up to 5 attempts.
    \end{enumerate}
    \item If both methods fail, the point is flagged as non-converged and assigned \texttt{NaN}.
\end{enumerate}

The same strategy is applied for all single-variable inversions:
\begin{itemize}
    \item $T(s, P)$: invert $s(P,T) = s_{\rm target}$,
    \item $P(\rho, T)$: invert $\rho(P,T) = \rho_{\rm target}$,
    \item $T(\rho, P)$: invert $\rho(P,T) = \rho_{\rm target}$ for $T$,
    \item $P(s, \rho)$: invert $s(P,T)$ and $\rho(P,T)$ for $P$,
    \item $T(s, \rho)$: invert similarly for $T$.
\end{itemize}

\subsubsection{Two-Dimensional Inversion}

For the simultaneous inversion $(s, \rho) \to (P, T)$, we solve a coupled system:
\begin{equation}
    \bm{F}(\log P, \log T) = \begin{pmatrix} s_{\rm VAL}(\log P, \log T)/s_{\rm target} - 1 \\ \log\rho_{\rm VAL}(\log P, \log T)/\log\rho_{\rm target} - 1 \end{pmatrix} = \bm{0},
\end{equation}
using \texttt{scipy.optimize.root} with the Powell hybrid method, or alternatively a Nelder--Mead simplex minimization of $|\bm{F}|^2$.

\subsection{Post-Inversion Corrections}\label{sec:glitch}

Numerical inversions may produce isolated outliers (``glitches'') or fail to converge at scattered points.
We apply a multi-stage correction pipeline along the composition ($Z$) axis at each $(s,P)$ or $(\rho,T)$ grid point:

\begin{enumerate}
    \item \textbf{Non-converged point interpolation:} Points where root-finding failed are replaced by linear (or quadratic) interpolation from the converged neighbors along the $Z$ axis.

    \item \textbf{Adaptive Hampel filter (two passes):} At each point $i$ along the $Z$ axis, we compute the local median $\tilde{y}$ and the median absolute deviation $\mathrm{MAD} = 1.4826\,\mathrm{median}(|y_j - \tilde{y}|)$ within a sliding window.
    Points satisfying $|y_i - \tilde{y}| > 3\,\mathrm{MAD}$ are flagged as outliers and replaced by the local median.
    Two passes are applied to catch cascaded glitches.

    \item \textbf{NaN interpolation:} Any remaining \texttt{NaN} values are filled by 1D linear interpolation (with extrapolation at the boundaries) using the valid neighbors.

    \item \textbf{Gaussian smoothing} (optional): For the coldest entropy slices ($s \leq 3\;k_B\;\mathrm{baryon}^{-1}$), a 1D Gaussian filter with $\sigma=3$ grid points is applied along the $Z$ axis to suppress residual noise from phase transitions in the water EOS.
\end{enumerate}

\subsection{Table Generation}

The complete inversion pipeline loops over all grid points of the target coordinate system.
For example, to generate $S$--$P$ tables, we loop over $(s_i, \log P_j)$ and, for each pair, solve for $\log T$ and then evaluate $\log\rho$ at all $(Y_k, Z_l)$ values.
The composition axis is swept sequentially, using the previous $Y$ solution as the initial guess for the next, while the $(s, P)$ axes are independent and can be parallelized.

%============================================================
\section{Derived Thermodynamic Quantities}\label{sec:derivatives}
%============================================================

All derivatives are computed by centered finite differences on the tabulated or inverted EOS.
We adopt the convention that $\ln$ denotes the natural logarithm and $\log$ denotes $\log_{10}$.

\subsection{Heat Capacities}

The specific heat at constant pressure is
\begin{equation}\label{eq:cp}
    C_P = T\left(\frac{\partial s}{\partial T}\right)_{\!P,Y,Z} = \frac{\partial s}{\partial \ln T}\bigg|_{P,Y,Z} \approx \frac{s(\log T + \delta) - s(\log T - \delta)}{2\,\delta\,\ln 10},
\end{equation}
where $\delta$ is a small step in $\log T$ (default $\delta = 0.1$).

Equivalently, when working in $(s,P)$ coordinates:
\begin{equation}
    C_P = \frac{2\,\Delta s}{\left[\log T(s+\Delta s) - \log T(s-\Delta s)\right]\,\ln 10},
\end{equation}
where $\Delta s$ is a step in entropy (default $\Delta s = 0.1\;k_B\;\mathrm{baryon}^{-1}$) and the unit conversion from $\mathrm{erg\;g^{-1}\;K^{-1}}$ to $k_B\;\mathrm{baryon}^{-1}$ is applied.

The specific heat at constant volume is computed analogously:
\begin{equation}\label{eq:cv}
    C_V = \frac{\partial s}{\partial \ln T}\bigg|_{\rho,Y,Z} \approx \frac{s(\rho, \log T + \delta) - s(\rho, \log T - \delta)}{2\,\delta\,\ln 10}.
\end{equation}

\subsection{Adiabatic Gradient}

The adiabatic temperature gradient is
\begin{equation}\label{eq:nabla_ad}
    \nabla_{\rm ad} \equiv \frac{\partial \log T}{\partial \log P}\bigg|_{s,Y,Z} \approx \frac{\log T(s, \log P + \delta) - \log T(s, \log P - \delta)}{2\,\delta},
\end{equation}
with default $\delta = 10^{-2}$.

The adiabatic entropy profile $s_{\rm ad}(P)$ is obtained by solving for the entropy that satisfies $\nabla = \nabla_{\rm ad}$ along a given pressure--temperature profile, using root-finding on the error function
\begin{equation}
    \varepsilon_{\nabla}(s) = \frac{\nabla_{\rm ad}(s, P)}{\nabla_{\rm profile}} - 1 = 0,
\end{equation}
where $\nabla_{\rm profile} = d\log T / d\log P$ is computed from the input $(P,T)$ profile.

\subsection{First Adiabatic Index}

The first adiabatic index is
\begin{equation}\label{eq:gamma1}
    \Gamma_1 \equiv \frac{\partial \ln P}{\partial \ln \rho}\bigg|_{s,Y,Z} = \frac{\partial \log P}{\partial \log \rho}\bigg|_{s,Y,Z} \approx \frac{2\,\delta}{\log\rho(s, \log P + \delta) - \log\rho(s, \log P - \delta)}.
\end{equation}

\subsection{Composition Derivatives and Convective Stability}\label{sec:composition}

The Schwarzschild criterion for convective instability requires the entropy derivative at constant $P$ and $T$:
\begin{equation}
    \frac{\partial s}{\partial Y}\bigg|_{P,T,Z} \approx \frac{s(Y + \delta_Y) - s(Y - \delta_Y)}{2\,\delta_Y}.
\end{equation}
An analogous expression holds for $\partial s/\partial Z|_{P,T,Y}$.

\subsubsection{Ledoux Derivatives via the Triple Product Rule}

The Ledoux criterion requires $\partial s/\partial Y|_{\rho,P,Z}$, which is not directly accessible from the tables.
We compute it using the triple product rule in the $(s,\rho)$ basis:
\begin{equation}\label{eq:triple}
    \frac{\partial s}{\partial Y}\bigg|_{\rho,P,Z} = -\frac{(\partial P/\partial Y)_{s,\rho,Z}}{(\partial P/\partial s)_{\rho,Y,Z}},
\end{equation}
where each partial derivative on the right-hand side is evaluated by centered finite differences on the $(s,\rho)$ tables.
Explicitly:
\begin{equation}
    \left(\frac{\partial P}{\partial s}\right)_{\!\rho,Y,Z} \approx \frac{P(s+\Delta s, \rho) - P(s-\Delta s, \rho)}{2\,\Delta s / (k_B/m_u)},
\end{equation}
\begin{equation}
    \left(\frac{\partial P}{\partial Y}\right)_{\!s,\rho,Z} \approx \frac{P(s, \rho, Y+\delta_Y) - P(s, \rho, Y-\delta_Y)}{2\,\delta_Y}.
\end{equation}
The same approach yields $\partial s/\partial Z|_{\rho,P,Y}$ by replacing $Y$ with $Z$ in the numerator.

\subsubsection{Susceptibility Coefficients}

The logarithmic susceptibility coefficients, evaluated in the $(\rho, T)$ basis, are
\begin{align}
    \chi_T &\equiv \frac{\partial \ln P}{\partial \ln T}\bigg|_{\rho,Y,Z}, \\
    \chi_Y &\equiv \frac{\partial \ln P}{\partial Y}\bigg|_{\rho,T,Z}, \\
    \chi_Z &\equiv \frac{\partial \ln P}{\partial Z}\bigg|_{\rho,T,Y}.
\end{align}
The temperature response to composition changes at constant $(\rho,P)$ is then
\begin{equation}
    \frac{\partial \ln T}{\partial Y}\bigg|_{\rho,P,Z} = \frac{\chi_Y}{\chi_T}, \qquad
    \frac{\partial \ln T}{\partial Z}\bigg|_{\rho,P,Y} = \frac{\chi_Z}{\chi_T}.
\end{equation}

\subsection{Density Response}

The Brunt--V\"ais\"al\"a frequency requires the density response to entropy changes at constant pressure:
\begin{equation}
    \frac{\partial \ln \rho}{\partial s}\bigg|_{P,Y,Z} \approx \frac{[\log\rho(s+\Delta s) - \log\rho(s-\Delta s)]\,\ln 10}{2\,\Delta s / (k_B/m_u)},
\end{equation}
as well as the thermal expansion coefficient
\begin{equation}
    \frac{\partial \log \rho}{\partial \log T}\bigg|_{P,Y,Z} \approx \frac{\log\rho(\log T + \delta) - \log\rho(\log T - \delta)}{2\,\delta}.
\end{equation}

\subsection{Chemical Potential and Energy Derivatives}

Thermodynamic consistency of the EOS can be verified via the fundamental relation.
In the $(s,\rho)$ basis:
\begin{equation}
    \left(\frac{\partial u}{\partial s}\right)_{\!\rho,Y} = T, \qquad
    -\left(\frac{\partial u}{\partial v}\right)_{\!s,Y} = P,
\end{equation}
where $v = 1/\rho$ is the specific volume.
The chemical potential differences are obtained from
\begin{equation}
    \left(\frac{\partial u}{\partial Y}\right)_{\!s,\rho,Z}, \qquad
    \left(\frac{\partial u}{\partial Z}\right)_{\!s,\rho,Y},
\end{equation}
each computed by centered finite differences.

%============================================================
\section{EOS Smoothing}\label{sec:smoothing}
%============================================================

The raw CMS19 and CD21 hydrogen--helium tables contain numerical artifacts from the bicubic spline interpolation used to stitch together the SCvH, QMD, and CP98 regimes.
These artifacts manifest as discontinuities in entropy and density gradients, particularly at the SCvH$\to$QMD boundary ($\rho \approx 0.05\;\mathrm{g\;cm^{-3}}$) and near the H$_2$ melting curve.
Similar artifacts exist in the water, methane, ammonia, and forsterite EOS tables near phase boundaries.

\subsection{Five-Step H--He Smoothing Pipeline}

For the H and He tables, we apply a five-step pipeline before constructing the interpolators:
\begin{enumerate}
    \item \textbf{Outlier isotherm replacement:} Isotherms with anomalously large deviations from their neighbors are identified and replaced by linear interpolation.
    \item \textbf{Low-pressure extrapolation:} Invalid cells at low pressures (below the solid/liquid boundary) are detected via interior gradient comparison and replaced by linear extrapolation from valid neighbors.
    \item \textbf{Low-pressure 2D Gaussian smoothing:} A 2D Gaussian filter is applied in the low-pressure region, blended with the original data via a tanh mask.
    \item \textbf{High-pressure extrapolation:} Invalid cells at high pressures (above the quantum correction limit) are similarly detected and extrapolated.
    \item \textbf{Targeted Gaussian smoothing:} A 2D Gaussian filter targets the dissociation/ionization zone using a gradient-adaptive mask.
\end{enumerate}
Each step operates on all three grids ($\log\rho$, $\log s$, $\log u$) simultaneously and is controlled by the \texttt{smooth\_hhe} flag in the constructor.

\subsection{Heavy-Element EOS Smoothing}

\subsubsection{AQUA Water EOS}
The AQUA water table exhibits discontinuities at low pressures and temperatures from the ice--vapor phase boundary.
A localized 2D Gaussian filter ($\sigma = 3$ grid points) is blended via a tanh mask concentrated in the low-$P$/low-$T$ corner:
\begin{equation}
    w_P = \tfrac{1}{2}\bigl[1 - \tanh\bigl((\log P - 8)/1.5\bigr)\bigr], \quad
    w_T = \tfrac{1}{2}\bigl[1 - \tanh\bigl((\log T - 3)/0.3\bigr)\bigr],
\end{equation}
with the combined mask $w = w_P \cdot w_T$ applied to blend the smoothed and original grids.

\subsubsection{CH$_4$ and NH$_3$ EOS}
The methane and ammonia tables exhibit negative or NaN entropy values at high pressures and low temperatures.
The smoothing procedure first interpolates invalid values from valid neighbors along each isotherm, then applies a targeted 2D Gaussian ($\sigma_T = 2$, $\sigma_P = 3$ grid points) blended via a tanh mask tuned to the discontinuity region.

\subsubsection{Mg$_2$SiO$_4$ (Forsterite) ANEOS}
The ANEOS forsterite table contains phase-transition discontinuities.
A gradient-adaptive approach detects discontinuities via the relative gradient magnitude, then applies a Gaussian filter selectively in those regions, combined with a low-pressure mask.

%============================================================
\section{Multi-Species Z EOS}\label{sec:zeos}
%============================================================

The \texttt{z\_eos} class loads heavy-element EOS tables directly from raw data files, supporting four species:
\begin{itemize}
    \item \textbf{Water} (AQUA): $P$--$T$ and $\rho$--$T$ bases, SI units converted to CGS.
    \item \textbf{Methane} (CH$_4$): $P$--$T$ basis, already in CGS.
    \item \textbf{Ammonia} (NH$_3$): $P$--$T$ basis, already in CGS.
    \item \textbf{Forsterite} (Mg$_2$SiO$_4$): ANEOS $P$--$T$ basis, $P$ in GPa converted to dyn/cm$^2$.
\end{itemize}
Each instance builds its own \texttt{RegularGridInterpolator} objects, allowing multiple independent \texttt{z\_eos} instances to coexist for metal-mixture calculations.
The axis ordering differs by species: water and forsterite use $(\log P, \log T)$, while CH$_4$ and NH$_3$ use $(\log T, \log P)$, handled transparently by an internal dispatch method.

%============================================================
\section{$P$--$T$ Dependent Hydrogen Molecular Weight}\label{sec:mu_h}
%============================================================

The CMS19 hydrogen EOS transitions from molecular H$_2$ ($\mu_{\rm H} = 2$) at low temperatures to atomic/ionized H ($\mu_{\rm H} = 1$) at high temperatures.
This transition affects the ideal entropy of mixing through the number fractions $x_i = f_i/\mu_i / \sum_j f_j/\mu_j$.

We determine the dissociation boundary empirically from the CMS19 table by locating the peak of the isobaric heat capacity $C_P$ along each isobar.
In the low-pressure regime ($\log P_{\rm cgs} \leq 10.5$), the dissociation peak is cleanly separated from the ionization peak, yielding the linear fit:
\begin{equation}\label{eq:dissoc}
    \log T_{\rm dissoc} = 2.966 + 0.097 \, \log P_{\rm cgs},
\end{equation}
with a maximum residual of 0.06\;dex over the fit range.
Below this boundary, $\mu_{\rm H} = 2$ (molecular H$_2$); above it, $\mu_{\rm H} = 1$ (atomic or ionized H).
The ionized species H$^+$ also uses $\mu_{\rm H} = 1$ since we count hydrogen atoms (not free electrons) as mixing particles.

When the \texttt{mu\_h\_vary} flag is enabled (default), the entropy of mixing is computed with the $P$--$T$ dependent $\mu_{\rm H}$.
Setting \texttt{mu\_h\_vary=False} recovers the legacy behavior ($\mu_{\rm H} = 1$ everywhere).

%============================================================
\section{Rhomboid-Adaptive Inversion Tables}\label{sec:rhomboid}
%============================================================

\subsection{The Rhomboid Problem}\label{sec:rhomboid_problem}

When inverting from the $P$--$T$ basis to alternative coordinate systems, the physical domain in the target coordinates is generally non-rectangular.
For example, in $(s, P)$ space, the entropy range at each pressure spans from $s_{\rm lo}(P) = s(P, T_{\rm max})$ to $s_{\rm hi}(P) = s(P, T_{\rm min})$, which varies by an order of magnitude across the pressure range.
A fixed rectangular $s$ grid wastes significant storage on unphysical cells (only $\sim$31\% valid for $(\rho, P)$ tables) and risks storing garbage values from failed inversions.

\subsection{Normalized Entropy Coordinate $\xi$}\label{sec:xi}

We address this by introducing a normalized coordinate $\xi \in [0, 1]$ that maps the physical domain onto a rectangle.
At each point along the non-varying axis (e.g., $\log P$ for S--P tables), the bounds of the varying axis are computed from the temperature boundaries of the $P$--$T$ table:
\begin{equation}\label{eq:xi}
    \xi = \frac{s - s_{\rm lo}(P, Y', Z)}{s_{\rm hi}(P, Y', Z) - s_{\rm lo}(P, Y', Z)},
\end{equation}
where $s_{\rm lo}$ and $s_{\rm hi}$ are stored as 3D arrays over $(\log P, Y', Z)$ and interpolated at query time.

The table is then built on a rectangular grid in $(\xi, \log P, Y', Z)$, and the \texttt{RegularGridInterpolator} operates on these rectangular axes.
To query the table with a physical entropy $s$, the caller converts $s \to \xi$ using the interpolated bounds, queries the RGI, and obtains $\log T(\xi, \log P, Y', Z)$.

\subsection{Application to Each Basis}\label{sec:xi_basis}

The $\xi$-mapping is applied to four inversion bases:

\begin{enumerate}
    \item \textbf{$S$--$P$ tables:} $\xi$ normalizes entropy at each $(\log P, Y', Z)$.
    Bounds: $s_{\rm lo}(P) = s(P, T_{\rm max})$, $s_{\rm hi}(P) = s(P, T_{\rm min})$.
    Stored: $\log T(\xi, \log P, Y', Z)$ in float32.

    \item \textbf{$\rho$--$T$ tables:} $\xi$ normalizes $\log\rho$ at each $(\log T, Y', Z)$.
    Bounds: $\rho_{\rm lo}(T) = \rho(P_{\rm min}, T)$, $\rho_{\rm hi}(T) = \rho(P_{\rm max}, T)$.
    Stored: $\log P(\xi, \log T, Y', Z)$.

    \item \textbf{$\rho$--$P$ tables:} $\xi$ normalizes $\log\rho$ at each $(\log P, Y', Z)$.
    Bounds: $\rho_{\rm lo}(P) = \rho(P, T_{\rm max})$, $\rho_{\rm hi}(P) = \rho(P, T_{\rm min})$.
    Stored: $\log T(\xi, \log P, Y', Z)$.

    \item \textbf{$S$--$\rho$ tables:} $\xi$ normalizes entropy at each $(\log\rho, Y', Z)$.
    Bounds computed by sweeping $T$ at each $\rho$ (via 1D $\rho$-inversion for $P$).
    Stored: $\log P(\xi, \log\rho, Y', Z)$ and $\log T(\xi, \log\rho, Y', Z)$.
    The 2D inversion uses bounded least-squares (Trust Region Reflective method).
\end{enumerate}

The bounds are per-$(Y', Z)$, not a global envelope, because the entropy and density ranges vary significantly with composition (e.g., $s_{\rm hi}$ at fixed $P$ varies by a factor of 4 between $Z=0$ and $Z=1$).

\subsection{NaN Repair and Memory Optimization}

After the root-finding loop, any remaining NaN cells (from convergence failures at the physical-domain boundary) are filled by a three-pass interpolation strategy:
\begin{enumerate}
    \item Along the $\xi$ axis (the most physical direction),
    \item Along the second coordinate axis ($\log P$, $\log T$, or $\log\rho$),
    \item Nearest-neighbor extrapolation for remaining edge cells.
\end{enumerate}
This guarantees zero NaNs in the final table.

All tables are stored in \texttt{float32} to halve memory.
For the S--P, $\rho$--T, and $\rho$--P tables, only the primary dependent variable ($\log T$ or $\log P$) is stored; all other quantities ($\log\rho$, $s$, $u$) are derived on-the-fly from the forward model at $(P, T)$.

\subsection{Warm-Start Solver Strategy}

During table generation, the solver uses a continuation strategy to accelerate convergence:
\begin{itemize}
    \item At the start of each new $Y'$ value, the guess cache is reset.
    \item For the first $Z$ step, Brent's method (guaranteed convergence with bracket $[\log T_{\rm min}, \log T_{\rm max}]$) establishes initial solutions.
    \item For subsequent $Z$ steps ($\Delta Z = 0.02$), Newton's method is warm-started from the previous $Z$ solution.  Since the composition perturbation is small, Newton typically converges in 2--4 iterations.
    \item If Newton fails or wanders outside the valid range, the solver falls back to Brent's method.
\end{itemize}

%============================================================
\section{P--T Basis Thermodynamic Identities}\label{sec:pt_derivs}
%============================================================

A key advance in the new framework is that \emph{all thermodynamic derivatives} are computed purely from $P$--$T$ basis quantities, avoiding the need for pre-computed inverted tables.
This is achieved through standard thermodynamic identities and the triple product rule.

\subsection{Base Quantities}\label{sec:base_derivs}

We define four dimensionless log-log derivatives, each evaluated by centered finite differences on the forward model at a given $(\log P, \log T, Y', Z)$:
\begin{align}
    a &\equiv \frac{\partial \log s}{\partial \log T}\bigg|_{P} = \frac{\log s(P, T{+}\delta T) - \log s(P, T{-}\delta T)}{2\,\delta\log T}, \label{eq:a} \\[4pt]
    b &\equiv \frac{\partial \log s}{\partial \log P}\bigg|_{T} = \frac{\log s(P{+}\delta P, T) - \log s(P{-}\delta P, T)}{2\,\delta\log P}, \label{eq:b} \\[4pt]
    c &\equiv \frac{\partial \log\rho}{\partial \log T}\bigg|_{P} = \frac{\log\rho(P, T{+}\delta T) - \log\rho(P, T{-}\delta T)}{2\,\delta\log T}, \label{eq:c} \\[4pt]
    d &\equiv \frac{\partial \log\rho}{\partial \log P}\bigg|_{T} = \frac{\log\rho(P{+}\delta P, T) - \log\rho(P{-}\delta P, T)}{2\,\delta\log P}, \label{eq:d}
\end{align}
with default step sizes $\delta\log T = \delta\log P = 10^{-2}$.

The composition derivatives (at constant $P, T$) are:
\begin{equation}
    S_Y \equiv \frac{\partial s}{\partial Y'}\bigg|_{P,T,Z}, \quad
    S_Z \equiv \frac{\partial s}{\partial Z}\bigg|_{P,T,Y'}, \quad
    \rho_Y \equiv \frac{\partial \log\rho}{\partial Y'}\bigg|_{P,T,Z}, \quad
    \rho_Z \equiv \frac{\partial \log\rho}{\partial Z}\bigg|_{P,T,Y'},
\end{equation}
each evaluated by centered finite differences with adaptive step sizing near the composition boundaries $[0, 1]$.

\subsection{Constrained Partial Derivative Identity}\label{sec:constrained_pde}

A recurring identity in the following derivations is the \emph{constrained partial derivative expansion}.
Suppose $f = f(x, y)$ and we wish to evaluate $(\partial f/\partial x)_g$ where $g = g(x, y)$ is held constant.
Along a path of constant $g$, $y$ is implicitly a function of $x$ via $g(x, y(x)) = \text{const}$, so
\begin{equation}\label{eq:constrained_pd}
    \left(\frac{\partial f}{\partial x}\right)_{\!g}
    = \left(\frac{\partial f}{\partial x}\right)_{\!y}
    + \left(\frac{\partial f}{\partial y}\right)_{\!x}
    \cdot \left(\frac{\partial y}{\partial x}\right)_{\!g},
\end{equation}
where the last factor follows from the triple product rule on $g(x,y)$:
\begin{equation}
    \left(\frac{\partial y}{\partial x}\right)_{\!g} = -\frac{(\partial g/\partial x)_y}{(\partial g/\partial y)_x}.
\end{equation}
This identity is used to derive $C_V$, the Ledoux stability derivatives, and the constrained composition derivatives below.

\subsection{Heat Capacities}\label{sec:heat_cap_id}

The isobaric heat capacity is
\begin{equation}\label{eq:cp_identity}
    C_P \equiv T\left(\frac{\partial s}{\partial T}\right)_{\!P} = \frac{\partial s}{\partial \ln T}\bigg|_P = s \cdot a.
\end{equation}
For the isochoric heat capacity $C_V = T(\partial s/\partial T)_\rho$, we apply equation~(\ref{eq:constrained_pd}) with $f = s$, $x = T$, $y = P$, and $g = \rho$:
\begin{equation}
    \left(\frac{\partial s}{\partial T}\right)_{\!\rho}
    = \left(\frac{\partial s}{\partial T}\right)_{\!P}
    + \left(\frac{\partial s}{\partial P}\right)_{\!T}
    \cdot \left(\frac{\partial P}{\partial T}\right)_{\!\rho}.
\end{equation}
The triple product rule on $\rho(P,T)$ gives
\begin{equation}
    \left(\frac{\partial P}{\partial T}\right)_\rho = -\frac{(\partial\rho/\partial T)_P}{(\partial\rho/\partial P)_T},
\end{equation}
which in log-derivative notation is $(\partial\log P/\partial\log T)_\rho = -c/d$.
Substituting:
\begin{equation}\label{eq:cv_identity}
    C_V = s \left(a - \frac{b\,c}{d}\right).
\end{equation}

\subsection{Adiabatic Gradient and First Adiabatic Index}

From the triple product rule on $s(P, T)$:
\begin{equation}\label{eq:nabla_ad_identity}
    \nabla_{\rm ad} \equiv \frac{\partial\log T}{\partial\log P}\bigg|_s = -\frac{b}{a}.
\end{equation}
The first adiabatic index follows from the chain rule for $(\partial P/\partial\rho)_s$:
\begin{equation}\label{eq:gamma1_identity}
    \Gamma_1 \equiv \frac{\partial\log P}{\partial\log\rho}\bigg|_s = \frac{1}{d - c\,b/a},
\end{equation}
which satisfies the identity $\Gamma_1 = \chi_\rho \cdot C_P / C_V$.

\subsection{Susceptibility Coefficients}

The coefficients $\chi_T$, $\chi_\rho$, $\chi_Y$, and $\chi_Z$ are obtained directly from the base quantities:
\begin{equation}\label{eq:chi_identity}
    \chi_T = -\frac{c}{d}, \quad
    \chi_\rho = \frac{1}{d}, \quad
    \chi_Y = -\frac{\rho_Y}{d}, \quad
    \chi_Z = -\frac{\rho_Z}{d}.
\end{equation}
The thermal expansion parameter is $\delta = -c$.

\subsection{Ledoux Stability Derivatives}\label{sec:ledoux_id}

The Ledoux criterion requires entropy derivatives at constant $(\rho, P)$, which constrain the temperature to adjust when composition changes.
At constant $P$ and $\rho$, a change $dY'$ induces a temperature change
\begin{equation}
    \frac{dT}{dY'}\bigg|_{\rho,P} = -\frac{(\partial\rho/\partial Y')_{P,T}}{(\partial\rho/\partial T)_P} = -T\,\ln 10\, \frac{\rho_Y}{c},
\end{equation}
where the factor of $\ln 10$ arises because $\rho_Y \equiv \partial(\log_{10}\rho)/\partial Y'$ is a log$_{10}$ derivative while $c$ is a log--log derivative (the $\ln 10$ in $\partial\rho/\partial Y' = \rho\,\ln 10\,\rho_Y$ does not cancel against $\partial\rho/\partial T = \rho\,c/T$ since $c$ is already dimensionless).
The Ledoux entropy derivative is then:
\begin{equation}\label{eq:ledoux_y}
    \frac{\partial s}{\partial Y'}\bigg|_{\rho,P,Z} = S_Y - s\,a\,\ln 10\,\frac{\rho_Y}{c},
\end{equation}
and analogously for $Z$:
\begin{equation}\label{eq:ledoux_z}
    \frac{\partial s}{\partial Z}\bigg|_{\rho,P,Y'} = S_Z - s\,a\,\ln 10\,\frac{\rho_Z}{c}.
\end{equation}

\subsection{Constrained Composition Derivatives}\label{sec:constrained}

The temperature and pressure responses to composition changes at constant $(s, \rho)$ are obtained from the implicit function theorem.
The two constraints $s(P, T, Y') = \text{const}$ and $\rho(P, T, Y') = \text{const}$ define a 2$\times$2 linear system in the actual (non-log) derivatives:
\begin{equation}
    \begin{pmatrix} s\,b/P & s\,a/T \\ \rho\,d/P & \rho\,c/T \end{pmatrix}
    \begin{pmatrix} dP \\ dT \end{pmatrix}
    = -\begin{pmatrix} S_Y \\ \rho\,\ln 10\,\rho_Y \end{pmatrix} dY',
\end{equation}
where the right-hand side uses $\partial\rho/\partial Y' = \rho\,\ln 10\,\rho_Y$ (since $\rho_Y$ is in $\log_{10}$).
Solving by Cramer's rule:
\begin{align}
    \frac{dT}{dY'}\bigg|_{s,\rho} &= \frac{T\,(d\,S_Y - s\,b\,\ln 10\,\rho_Y)}{s\,(b\,c - a\,d)}, \label{eq:dTdY_srho} \\[4pt]
    \frac{dP}{dY'}\bigg|_{s,\rho} &= \frac{P\,(s\,a\,\ln 10\,\rho_Y - c\,S_Y)}{s\,(b\,c - a\,d)}. \label{eq:dPdY_srho}
\end{align}
The chemical potential derivatives follow from the chain rule:
\begin{equation}\label{eq:dUdY_srho}
    \frac{\partial u}{\partial Y'}\bigg|_{s,\rho} = U_P \cdot \frac{d\log P}{dY'}\bigg|_{s,\rho} + U_T \cdot \frac{d\log T}{dY'}\bigg|_{s,\rho} + U_Y,
\end{equation}
where $U_P = (\partial u/\partial\log P)_T$, $U_T = (\partial u/\partial\log T)_P$, $U_Y = (\partial u/\partial Y')_{P,T}$, and
$d\log P/dY' = (dP/dY')/(P\,\ln 10)$, $d\log T/dY' = (dT/dY')/(T\,\ln 10)$.
Analogous expressions hold with $Z$ replacing $Y'$.

\subsection{Smoothed Derivatives}\label{sec:savgol}

The residual kinks in the CMS19/CD21 tables (particularly at the SCvH$\to$QMD stitching boundary) can produce noisy finite differences.
To mitigate this, we optionally replace the 2-point central difference with a 5-point Savitzky--Golay derivative filter.

For a derivative $df/dx$ at point $x_0$, we evaluate $f$ on a 5-point stencil $\{x_0 - 2h, x_0 - h, x_0, x_0 + h, x_0 + 2h\}$ and compute
\begin{equation}\label{eq:savgol}
    \frac{df}{dx}\bigg|_{x_0} = \frac{-2f_{-2} - f_{-1} + f_1 + 2f_2}{10\,h},
\end{equation}
which is the first-derivative coefficient of a degree-2 polynomial fit to the five points (equivalent to a weighted moving average of slopes).
This has the same formal order of accuracy as the 2-point central difference but suppresses high-frequency noise.

\subsection{Summary of Derivative Methods}

Table~\ref{tab:derivatives} summarizes all 18 thermodynamic derivatives, their finite-difference basis, the inversion required, and their physical role.

\begin{deluxetable*}{llcll}
\tablecaption{Thermodynamic Derivatives and Their Finite-Difference Bases\label{tab:derivatives}}
\tablecolumns{5}
\tablewidth{0pt}
\tablehead{
    \colhead{Method} &
    \colhead{Symbol} &
    \colhead{FD Basis} &
    \colhead{Inversion Required} &
    \colhead{Description}
}
\startdata
\multicolumn{5}{c}{\textbf{Pressure--Temperature Basis}} \\
\hline
\texttt{get\_cp}        & $C_P = \partial S / \partial \ln T |_P$                   & $P$--$T$ & None                       & Isobaric specific heat \\
\texttt{get\_delta}     & $\delta = -\partial \log\rho / \partial \log T |_P$       & $P$--$T$ & None                       & Thermal expansion parameter \\
\texttt{get\_dsdy\_pt}  & $\partial S / \partial Y' |_{P,T}$                        & $P$--$T$ & None                       & Schwarzschild stability (Y) \\
\texttt{get\_dsdz\_pt}  & $\partial S / \partial Z |_{P,T}$                         & $P$--$T$ & None                       & Schwarzschild stability (Z) \\
\hline
\multicolumn{5}{c}{\textbf{Entropy--Pressure Basis}} \\
\hline
\texttt{get\_nabla\_ad} & $\nabla_{\rm ad} = \partial \log T / \partial \log P |_S$ & $S$--$P$ & \texttt{get\_logt\_sp}     & Adiabatic temperature gradient \\
\texttt{get\_gamma1}    & $\Gamma_1 = \partial \log P / \partial \log\rho |_S$      & $S$--$P$ & \texttt{get\_logt\_sp}     & First adiabatic index \\
\texttt{get\_dtds\_sp}  & $\partial T / \partial S |_P$                             & $S$--$P$ & \texttt{get\_logt\_sp}     & Temperature--entropy response \\
\hline
\multicolumn{5}{c}{\textbf{Density--Temperature Basis}} \\
\hline
\texttt{get\_cv}        & $C_V = \partial S / \partial \ln T |_\rho$                & $\rho$--$T$ & \texttt{get\_logp\_rhot}   & Isochoric specific heat \\
\texttt{get\_chi\_T}    & $\chi_T = \partial \log P / \partial \log T |_\rho$       & $\rho$--$T$ & \texttt{get\_logp\_rhot}   & Thermal susceptibility \\
\texttt{get\_chi\_rho}  & $\chi_\rho = \partial \log P / \partial \log\rho |_T$     & $\rho$--$T$ & \texttt{get\_logp\_rhot}   & Density susceptibility \\
\texttt{get\_chi\_Y}    & $\chi_Y = \partial \log P / \partial Y' |_{\rho,T}$       & $\rho$--$T$ & \texttt{get\_logp\_rhot}   & Helium susceptibility \\
\texttt{get\_chi\_Z}    & $\chi_Z = \partial \log P / \partial Z |_{\rho,T}$        & $\rho$--$T$ & \texttt{get\_logp\_rhot}   & Metal susceptibility \\
\hline
\multicolumn{5}{c}{\textbf{Density--Pressure Basis}} \\
\hline
\texttt{get\_dsdy\_rhop} & $\partial S / \partial Y' |_{\rho,P}$                    & $\rho$--$P$ & Constrained $\rho(P,T)$   & Ledoux stability (Y) \\
\texttt{get\_dsdz\_rhop} & $\partial S / \partial Z |_{\rho,P}$                     & $\rho$--$P$ & Constrained $\rho(P,T)$   & Ledoux stability (Z) \\
\hline
\multicolumn{5}{c}{\textbf{Entropy--Density Basis}} \\
\hline
\texttt{get\_dtdy\_srho} & $\partial T / \partial Y' |_{S,\rho}$                    & $S$--$\rho$ & \texttt{get\_logp\_logt\_srho} & Temperature expansion (Y) \\
\texttt{get\_dtdz\_srho} & $\partial T / \partial Z |_{S,\rho}$                     & $S$--$\rho$ & \texttt{get\_logp\_logt\_srho} & Temperature expansion (Z) \\
\texttt{get\_dudy\_srho} & $\partial U / \partial Y' |_{S,\rho}$                    & $S$--$\rho$ & \texttt{get\_logp\_logt\_srho} & Chemical potential (Y) \\
\texttt{get\_dudz\_srho} & $\partial U / \partial Z |_{S,\rho}$                     & $S$--$\rho$ & \texttt{get\_logp\_logt\_srho} & Chemical potential (Z) \\
\enddata
\tablecomments{
All derivatives accept \texttt{method='finite\_difference'} (default) or \texttt{method='identity'}.
The finite-difference method uses the indicated inversion basis with a 5-point Savitzky--Golay stencil (\texttt{smooth=True}, default) or a 2-point central difference (\texttt{smooth=False}).
The identity method computes all quantities from the $P$--$T$ basis using thermodynamic identities (\S\ref{sec:pt_derivs}), avoiding inversions entirely.
Derivatives in the $\rho$--$P$ basis use a constrained finite difference: at constant $P$ and $\rho$, $T$ is adjusted via 1-D root-finding on $\rho(P,T) = \mathrm{const}$.
Derivatives in the $S$--$\rho$ basis use 2-D least-squares inversion via \texttt{scipy.optimize.least\_squares} with the Trust Region Reflective method.
}
\end{deluxetable*}


%============================================================
\section{Thermodynamic Consistency Verification}\label{sec:thermo_consistency}
%============================================================

Because the CMS19/CD21 hydrogen--helium tables and the AQUA water table are constructed from independent ab~initio calculations, the three independently tabulated quantities---specific entropy $s(P,T)$, density $\rho(P,T)$, and specific internal energy $u(P,T)$---are not guaranteed to satisfy the first law of thermodynamics exactly.
We quantify these inconsistencies with three tests derived from the fundamental relation $du = T\,ds - P\,dV$, where $V = 1/\rho$ is the specific volume.

\subsection{Derivation of Consistency Relations}

The first law of thermodynamics for a single-component fluid is
\begin{equation}\label{eq:first_law}
    du = T\,ds - P\,dV = T\,ds + \frac{P}{\rho^2}\,d\rho\,,
\end{equation}
where the second equality uses $dV = d(1/\rho) = -\rho^{-2}\,d\rho$.

\paragraph{Test A: First law at constant pressure.}
Restricting Equation~\eqref{eq:first_law} to constant $P$:
\begin{equation}\label{eq:first_law_constP}
    \left.\frac{\partial u}{\partial T}\right|_P
    = T \left.\frac{\partial s}{\partial T}\right|_P
    + \frac{P}{\rho^2} \left.\frac{\partial \rho}{\partial T}\right|_P \,.
\end{equation}
We convert to derivatives with respect to $\log_{10} T$ by multiplying both sides by $dT/d\!\log_{10}\!T = T \ln 10$:
\begin{equation}\label{eq:test_A}
    \boxed{
    \left.\frac{\partial u}{\partial \log T}\right|_P
    = T \left.\frac{\partial s}{\partial \log T}\right|_P
    + \frac{P}{\rho^2} \left.\frac{\partial \rho}{\partial \log T}\right|_P
    }
\end{equation}
where all derivatives are of the \emph{linear} (not logarithmic) quantities $s$, $\rho$, $u$ with respect to $\log_{10} T$ at constant $\log_{10} P$.

\paragraph{Test B: First law at constant temperature.}
The analogous relation at constant $T$, with derivatives in $\log_{10} P$:
\begin{equation}\label{eq:test_B}
    \boxed{
    \left.\frac{\partial u}{\partial \log P}\right|_T
    = T \left.\frac{\partial s}{\partial \log P}\right|_T
    + \frac{P}{\rho^2} \left.\frac{\partial \rho}{\partial \log P}\right|_T
    }
\end{equation}

\paragraph{Test C: Maxwell relation from the Gibbs free energy.}
The Gibbs free energy satisfies $dG = V\,dP - s\,dT$, so the equality of cross-derivatives gives
\begin{equation}
    \left.\frac{\partial V}{\partial T}\right|_P
    = -\left.\frac{\partial s}{\partial P}\right|_T \,.
\end{equation}
Substituting $V = 1/\rho$ so that $(\partial V/\partial T)|_P = -\rho^{-2}(\partial\rho/\partial T)|_P$, and converting to $\log_{10}$ derivatives (multiply both sides by $TP\ln 10$):
\begin{equation}\label{eq:test_C}
    \boxed{
    \frac{P}{\rho^2} \left.\frac{\partial \rho}{\partial \log T}\right|_P
    = T \left.\frac{\partial s}{\partial \log P}\right|_T
    }
\end{equation}

\subsection{Numerical Implementation}

All derivatives are evaluated via centered finite differences with step $h = 0.01$~dex in $\log_{10}$ space:
\begin{equation}
    \left.\frac{\partial f}{\partial \log T}\right|_P
    \approx \frac{f(\log P,\, \log T + h) - f(\log P,\, \log T - h)}{2h}\,,
\end{equation}
where $f$ is $s$, $\rho$, or $u$ in \emph{linear} units (erg\,g$^{-1}$\,K$^{-1}$, g\,cm$^{-3}$, erg\,g$^{-1}$, respectively).
The step size $h = 0.01$ corresponds to a $\sim\!2.3\%$ change in the physical variable and matches the finite-difference convention used in the derivative routines of Section~\ref{sec:derivatives}.

For each test, we define the left-hand side (LHS) and right-hand side (RHS) of the boxed equations above, and compute a symmetric fractional error:
\begin{equation}\label{eq:frac_error}
    \varepsilon = \frac{|\text{LHS} - \text{RHS}|}
    {\tfrac{1}{2}(|\text{LHS}| + |\text{RHS}|) + \epsilon}\,,
\end{equation}
where $\epsilon = 10^{-30}$ prevents division by zero.

Tests are evaluated on a uniform grid in $(\log P,\, \log T)$ spanning $\log P \in [6, 14]$ (dyn\,cm$^{-2}$) with step 0.05~dex, and $\log T \in [2.5, 5.5]$ (K) with step 0.025~dex, yielding $161 \times 121 = 19{,}481$ test points per composition.
Grid points where any of the six stencil evaluations (center $\pm\, h$ in both $\log T$ and $\log P$) return NaN are excluded.

\subsection{Results}

We evaluate the three consistency tests for H--He mixtures at $Y' = 0.275$ and $Z = 0$ using both the CMS19 and CD21 hydrogen tables, with no smoothing applied to the raw tables.  We also test H--He--water (AQUA) mixtures at $Z = 0.017$ and $Z = 0.051$ to assess the impact of the volume addition law mixing with heavy elements.

For the pure H--He case (CD21, $Y' = 0.275$, $Z = 0$), the median fractional errors are:
\begin{itemize}
    \item Test A (first law, const $P$): $\tilde{\varepsilon} \approx 1.0\%$, $\varepsilon_{95} \approx 25\%$
    \item Test B (first law, const $T$): $\tilde{\varepsilon} \approx 2.5\%$, $\varepsilon_{95} \approx 146\%$
    \item Test C (Maxwell relation): $\tilde{\varepsilon} \approx 1.1\%$, $\varepsilon_{95} \approx 46\%$
\end{itemize}
The largest violations are concentrated along the H$_2$ dissociation boundary ($\log T \approx 2.97 + 0.097 \log P_{\rm cgs}$), at table grid-node kinks (notably $\log T = 5.0$, where the bilinear interpolation preserves a slope discontinuity in the underlying table), and at the high-pressure domain edges where the tables run out of valid data.

Adding metals via the volume addition law slightly degrades the constant-$T$ test (Test~B) because the independently tabulated AQUA water energy introduces additional inconsistency along the pressure axis, while the constant-$P$ test (Test~A) and Maxwell relation (Test~C) remain largely unchanged.

%============================================================
\section{Ideal Gas Equation of State}\label{sec:ideal_eos}
%============================================================

The ideal gas EOS provides analytic thermodynamic functions for use as initial guesses in the numerical inversions described in \S\ref{sec:inversion}, and as a reference limit for validating the non-ideal tables.
Two classes are implemented: a monatomic ideal gas based on the Sackur--Tetrode equation (\texttt{IdealEOS}), and a molecular extension that includes rotational and vibrational degrees of freedom (\texttt{MolecularIdealEOS}).

\subsection{Monatomic Ideal Gas (Sackur--Tetrode)}\label{sec:sackur_tetrode}

For a monatomic ideal gas of molecular weight $\mu$ (in atomic mass units), the translational entropy per unit mass is derived from the Sackur--Tetrode equation.  In the $(P,T)$ basis, the entropy per particle in units of $k_B$ is
\begin{equation}\label{eq:sackur_tetrode}
    \frac{s_{\rm trans}^{(p)}}{k_B} = C_0^{(PT)} + \tfrac{5}{2}\ln T - \ln P + \tfrac{3}{2}\ln\mu,
\end{equation}
where $T$ is in kelvin, $P$ in $\mathrm{dyn\;cm^{-2}}$, $\mu$ is the dimensionless molecular weight, and
\begin{equation}\label{eq:c0pt}
    C_0^{(PT)} = \tfrac{5}{2} + \tfrac{5}{2}\ln k_B + \tfrac{3}{2}\ln(2\pi) + \tfrac{3}{2}\ln m_u - 3\ln h,
\end{equation}
with $k_B = 1.380649\times10^{-16}\;\mathrm{erg\;K^{-1}}$ (Boltzmann constant), $m_u = 1.66054\times10^{-24}\;\mathrm{g}$ (atomic mass unit), and $h = 6.62607\times10^{-27}\;\mathrm{erg\;s}$ (Planck constant; all exact in the 2019 SI).
The specific entropy in $\mathrm{erg\;g^{-1}\;K^{-1}}$ is then
\begin{equation}
    s_{\rm trans} = \frac{s_{\rm trans}^{(p)}}{\mu\, m_u / k_B}.
\end{equation}

Equivalent expressions in the $(\rho,T)$ and $(\rho,P)$ bases are obtained by substituting the ideal gas law $P = \rho\, k_B\, T / (\mu\, m_u)$:
\begin{align}
    \frac{s_{\rm trans}^{(p)}}{k_B}\bigg|_{\rho T} &= C_0^{(\rho T)} + \tfrac{3}{2}\ln T - \ln\rho + \tfrac{5}{2}\ln\mu, \label{eq:s_rhot} \\
    \frac{s_{\rm trans}^{(p)}}{k_B}\bigg|_{\rho P} &= C_0^{(\rho P)} + \tfrac{3}{2}\ln P + 4\ln\mu - \tfrac{5}{2}\ln\rho, \label{eq:s_rhop}
\end{align}
where $C_0^{(\rho T)} = C_0^{(PT)} + \ln m_u - \ln k_B$ and $C_0^{(\rho P)} = C_0^{(PT)} + \tfrac{5}{2}\ln m_u - \tfrac{5}{2}\ln k_B$.
The three constants are derived from a single base constant $C_0^{(PT)}$ via the ideal gas law, which guarantees mutual thermodynamic consistency.

The remaining thermodynamic quantities follow from standard ideal gas relations:
\begin{align}
    \rho &= P\,\mu\,m_u / (k_B T), \\
    u &= \tfrac{3}{2}\,(k_B / m_u)\, T / \mu, \\
    c_v &= \tfrac{3}{2}\, R_{\rm id} / \mu, \qquad c_p = \tfrac{5}{2}\, R_{\rm id} / \mu, \\
    \nabla_{\rm ad} &\equiv \frac{\partial \ln T}{\partial \ln P}\bigg|_s = \frac{\gamma-1}{\gamma} = \frac{2}{5},
\end{align}
where $R_{\rm id} \equiv k_B / m_u \approx 8.314\times10^{7}\;\mathrm{erg\;g^{-1}\;K^{-1}}$ is the gas constant per gram of unit-molecular-weight gas, and $\gamma = c_p/c_v = 5/3$.

The analytic inversions $T(s,P)$, $T(s,\rho)$, $P(s,\rho)$, and $\rho(s,P)$ are obtained by algebraic rearrangement of Equations~\eqref{eq:sackur_tetrode}--\eqref{eq:s_rhop}.
All inversions accept entropy in units of $k_B$ per baryon ($k_B / m_u$), consistent with the convention used by the non-ideal EOS modules.

\subsection{Molecular Ideal Gas: Rotational and Vibrational Contributions}\label{sec:molecular_eos}

For molecular species, the partition function factorizes as $q = q_{\rm trans}\, q_{\rm rot}\, q_{\rm vib}$, so the thermodynamic potentials are additive:
\begin{equation}
    s = s_{\rm trans} + s_{\rm rot} + s_{\rm vib}, \qquad u = u_{\rm trans} + u_{\rm rot} + u_{\rm vib}.
\end{equation}
We implement this as a subclass \texttt{MolecularIdealEOS} that inherits the translational (Sackur--Tetrode) contributions from \texttt{IdealEOS} and adds rotational and vibrational terms.
The ideal gas law relating $P$, $\rho$, and $T$ is unchanged by internal degrees of freedom.

\subsubsection{Rotational Contributions (Classical Limit)}\label{sec:rot}

In the classical high-temperature limit ($T \gg \Theta_{\rm rot}$), the rotational entropy per particle is
\begin{equation}\label{eq:s_rot_lin}
    \frac{s_{\rm rot}}{k_B} = \ln\frac{T}{\sigma\,\Theta_{\rm rot}} + 1 \qquad (\text{linear molecule}),
\end{equation}
\begin{equation}\label{eq:s_rot_nonlin}
    \frac{s_{\rm rot}}{k_B} = \ln\left[\frac{\sqrt{\pi}}{\sigma}\left(\frac{T^3}{\Theta_A\,\Theta_B\,\Theta_C}\right)^{1/2}\right] + \frac{3}{2} \qquad (\text{nonlinear}),
\end{equation}
where $\sigma$ is the rotational symmetry number and $\Theta_{\rm rot}$, $\Theta_A$, $\Theta_B$, $\Theta_C$ are the characteristic rotational temperatures defined by $\Theta = hcB/k_B$ for rotational constant $B$.
The corresponding heat capacities per particle are $c_{v,\rm rot} = k_B$ (linear; 2 rotational degrees of freedom) and $c_{v,\rm rot} = \frac{3}{2}k_B$ (nonlinear; 3 degrees of freedom).
Monatomic species have zero rotational contribution.

The classical limit is valid for all temperatures of interest in planetary interiors.
For H$_2$, $\Theta_{\rm rot} = 87.6\;\mathrm{K}$ and the quantum regime only matters below $\sim 100\;\mathrm{K}$; for heavier molecules $\Theta_{\rm rot} \ll 1\;\mathrm{K}$.

\subsubsection{Vibrational Contributions (Quantum Harmonic Oscillator)}\label{sec:vib}

Each vibrational mode $i$ with characteristic temperature $\Theta_{v,i}$ contributes according to the quantum harmonic oscillator (Einstein model).
Defining $x_i \equiv \Theta_{v,i}/T$, the per-particle contributions are
\begin{align}
    \frac{s_{{\rm vib},i}}{k_B} &= \frac{x_i}{e^{x_i}-1} - \ln\!\left(1 - e^{-x_i}\right), \label{eq:s_vib} \\
    \frac{u_{{\rm vib},i}}{k_B} &= \frac{\Theta_{v,i}}{e^{x_i}-1}, \label{eq:u_vib} \\
    \frac{c_{v,{\rm vib},i}}{k_B} &= \frac{x_i^2\, e^{x_i}}{(e^{x_i}-1)^2}. \label{eq:cv_vib}
\end{align}
A molecule with $N$ atoms has $3N-5$ vibrational modes (linear) or $3N-6$ modes (nonlinear).
The total vibrational contribution is the sum over all modes.

In the high-temperature limit ($T \gg \Theta_{v,i}$, i.e.\ $x_i \to 0$), each mode contributes $c_{v,{\rm vib},i} \to k_B$ per particle, recovering the classical equipartition value.
In the low-temperature limit ($T \ll \Theta_{v,i}$, i.e.\ $x_i \to \infty$), the vibrational mode freezes out: $c_{v,{\rm vib},i} \to 0$.
This temperature-dependent freeze-out means that the heat capacity and adiabatic gradient are no longer constant but vary with temperature:
\begin{equation}\label{eq:grad_ad_mol}
    \nabla_{\rm ad} = 1 - \frac{c_v}{c_p} = \frac{R_{\rm id}/\mu}{c_v + R_{\rm id}/\mu},
\end{equation}
where $c_v = c_{v,\rm trans} + c_{v,\rm rot} + c_{v,\rm vib}$ and $c_p = c_v + R_{\rm id}/\mu$ (the ideal gas relation $c_p - c_v = R/\mu$ holds regardless of internal degrees of freedom).

\subsubsection{Characteristic Temperatures}\label{sec:theta}

The characteristic temperatures are computed from spectroscopic constants tabulated in the NIST Computational Chemistry Comparison and Benchmark Database (CCCBDB).
We use \emph{harmonic} vibrational frequencies $\omega_e$ rather than fundamental (anharmonic) frequencies $\nu_0$, since the harmonic oscillator partition function (Equations~\ref{eq:s_vib}--\ref{eq:cv_vib}) assumes equally spaced energy levels; pairing it with anharmonic frequencies would be internally inconsistent.
The conversion is $\Theta_v = hc\,\omega_e / k_B$, with $hc/k_B = 1.4388\;\mathrm{cm\;K}$.

Table~\ref{tab:species} lists the predefined molecular species and their spectroscopic parameters.

\begin{deluxetable}{lccccc}
\tablecaption{Molecular species and their spectroscopic parameters from the NIST CCCBDB.  Rotational constants $B$ (or $A$, $B$, $C$ for nonlinear molecules) and harmonic vibrational frequencies $\omega_e$ are converted to characteristic temperatures $\Theta$ via $\Theta = hc\,B/k_B$ and $\Theta_v = hc\,\omega_e/k_B$.\label{tab:species}}
\tablehead{
    \colhead{Species} & \colhead{$\mu$} & \colhead{Geometry} & \colhead{$\sigma$} & \colhead{$\Theta_{\rm rot}$ (K)} & \colhead{$\Theta_{v}$ (K)}
}
\startdata
H$_2$    & 2.016  & linear    & 2 & 87.6              & 6332 \\
He       & 4.003  & monatomic & -- & --                & -- \\
H$_2$O   & 18.015 & nonlinear & 2 & 40.1, 20.9, 13.4 & 2373, 5514, 5674 \\
Fe       & 55.845 & monatomic & -- & --                & -- \\
MgSiO$_3$ & 100.39 & monatomic & -- & --               & -- \\
\enddata
\tablecomments{The H$_2$ rotational constant $B = 60.853\;\mathrm{cm^{-1}}$ and harmonic frequency $\omega_e = 4401\;\mathrm{cm^{-1}}$.  The H$_2$O rotational constants are $A = 27.877$, $B = 14.512$, $C = 9.285\;\mathrm{cm^{-1}}$; harmonic frequencies are $\omega_1 = 1649$, $\omega_2 = 3832$, $\omega_3 = 3943\;\mathrm{cm^{-1}}$.  Monatomic species (He, Fe, MgSiO$_3$) have no rotational or vibrational contributions and reduce identically to the Sackur--Tetrode result.}
\end{deluxetable}

\subsubsection{Inversions}\label{sec:mol_inversions}

With vibrational contributions, the entropy $s(P,T)$ is no longer analytically invertible for $T$ because the Einstein function (Equation~\ref{eq:s_vib}) is transcendental in $T$.
We solve for $T$ via Newton--Raphson iteration, using the monatomic Sackur--Tetrode analytic inversion as the initial guess.
The Jacobian
\begin{equation}
    \frac{\partial s}{\partial \log T}\bigg|_P = \left(\tfrac{5}{2} + \tfrac{n_{\rm rot}}{2}\right) \frac{k_B}{\mu\, m_u}\,\ln 10 + c_{v,\rm vib}(T)\,\ln 10
\end{equation}
is available analytically, where $n_{\rm rot}$ is the number of rotational degrees of freedom (0, 2, or 3).
For monatomic species the Newton loop performs zero iterations since the Sackur--Tetrode guess is exact, so there is no performance penalty.
For molecular species, convergence to a fractional tolerance of $10^{-12}$ typically requires 2--5 iterations.

The $T(s,\rho)$, $P(s,\rho)$, and $\rho(s,P)$ inversions are handled similarly: Newton iteration on $s(\rho,T)$ yields $T$, from which the remaining variable follows via the ideal gas law.

\subsection{Ideal H--He Mixture}\label{sec:ideal_hhe}

The \texttt{IdealHHeMix} class constructs a binary ideal gas mixture of hydrogen ($\mu_{\rm H} = 2.016$) and helium ($\mu_{\rm He} = 4.003$) using mass-fraction--weighted specific quantities:
\begin{equation}
    s_{\rm mix} = (1-Y')\,s_{\rm H}(P,T) + Y'\,s_{\rm He}(P,T) + s_{\rm mix,id},
\end{equation}
where $s_{\rm H}$ and $s_{\rm He}$ are the pure-species entropies (from \texttt{IdealEOS} by default) and $s_{\rm mix,id}$ is the ideal entropy of mixing (Equation~\ref{eq:smix_id_y}).
The mixture density follows the volume addition law:
\begin{equation}
    \frac{1}{\rho_{\rm mix}} = \frac{1-Y'}{\rho_{\rm H}(P,T)} + \frac{Y'}{\rho_{\rm He}(P,T)}.
\end{equation}
Internal energy is mixed linearly: $u_{\rm mix} = (1-Y')\,u_{\rm H} + Y'\,u_{\rm He}$.

The constructor accepts optional \texttt{eos\_h} and \texttt{eos\_he} arguments, allowing the hydrogen component to be replaced with a \texttt{MolecularIdealEOS} instance (e.g., molecular H$_2$ with rotational and vibrational modes) when a more physical ideal gas reference is desired.

\subsection{Verification}\label{sec:ideal_verification}

Both \texttt{IdealEOS} and \texttt{MolecularIdealEOS} are verified by a suite of self-consistency tests:

\begin{enumerate}
    \item \textbf{Entropy self-consistency:} $s(P,T) = s(\rho,T) = s(\rho,P)$ for thermodynamically consistent $(P,T,\rho)$ triples, verified to $\lesssim 10^{-15}$ fractional error.
    \item \textbf{Round-trip inversions:} $T \to s(P,T) \to T(s,P)$ recovers the original $T$ to $\lesssim 10^{-10}$ in $\log_{10} T$; similarly for all other inversion paths.
    \item \textbf{First law of thermodynamics:} At constant $P$, the centered finite-difference relation
    $\partial u / \partial \log T |_P = T\,\partial s / \partial \log T |_P + (P/\rho^2)\,\partial\rho/\partial\log T |_P$
    holds to $\lesssim 10^{-10}$ fractional error; similarly at constant $T$.
    \item \textbf{Maxwell relation:} $(P/\rho^2)\,\partial\rho/\partial\log T|_P = T\,\partial s/\partial\log P|_T$, verified to $\lesssim 10^{-10}$.
    \item \textbf{Classical equipartition limits:} At $T \gg \max(\Theta_{v,i})$, the heat capacity approaches the classical value---$c_v \to \frac{7}{2}\,R_{\rm id}/\mu$ for H$_2$ (3 translational + 2 rotational + 2 vibrational degrees of freedom) and $c_v \to 6\,R_{\rm id}/\mu$ for H$_2$O (3 + 3 + 6 degrees of freedom)---to better than 0.01\%.
    \item \textbf{Vibrational freeze-out:} At $T \ll \min(\Theta_{v,i})$, the heat capacity reduces to the translational + rotational value: $c_v \to \frac{5}{2}\,R_{\rm id}/\mu$ for H$_2$ and $c_v \to 3\,R_{\rm id}/\mu$ for H$_2$O.
    \item \textbf{NIST comparison:} The H$_2$ heat capacity at 300~K and 1~atm, $c_p = 29.1\;\mathrm{J\;mol^{-1}\;K^{-1}}$, agrees with the NIST tabulated value of $28.85\;\mathrm{J\;mol^{-1}\;K^{-1}}$ to $\lesssim 1\%$.  The residual difference is expected from the harmonic approximation.
    \item \textbf{Adiabatic gradient:} $\nabla_{\rm ad}$ transitions smoothly from 2/7 (H$_2$, low $T$) to 2/9 (H$_2$, high $T$) as vibrational modes activate, matching the analytic prediction $\nabla_{\rm ad} = 2/(2 + f)$ where $f$ is the number of active quadratic degrees of freedom.
\end{enumerate}

%============================================================
\section{Summary}\label{sec:summary}
%============================================================

We have described a comprehensive EOS framework for multi-component planetary mixtures.
The key advances over previous implementations are:
\begin{enumerate}
    \item \textbf{Targeted smoothing} of the H--He, water, CH$_4$, NH$_3$, and Mg$_2$SiO$_4$ tables to remove numerical artifacts from phase-boundary stitching while preserving physical features.
    \item \textbf{Multi-species metal mixing} via the volume addition law with independent \texttt{z\_eos} instances for water, methane, ammonia, and forsterite.
    \item \textbf{$P$--$T$ dependent hydrogen molecular weight} $\mu_{\rm H}(P,T)$ for the ideal entropy of mixing, switching from $\mu_{\rm H} = 2$ (molecular H$_2$) to $\mu_{\rm H} = 1$ (atomic/ionized H) at the empirically determined dissociation boundary.
    \item \textbf{Rhomboid-adaptive inversion tables} using a normalized coordinate $\xi \in [0,1]$ that maps the non-rectangular physical domain onto a rectangle suitable for \texttt{RegularGridInterpolator}.  This eliminates wasted storage on unphysical cells and avoids failed inversions at domain boundaries.
    \item \textbf{All thermodynamic derivatives from the $P$--$T$ basis}, using thermodynamic identities rather than inverted tables.  This makes derivative computation fast, independent of table availability, and avoids compounding interpolation errors through multiple inversion stages.
    \item \textbf{Savitzky--Golay smoothed finite differences} to suppress derivative noise from residual EOS table artifacts.
    \item \textbf{Molecular ideal gas EOS} with rotational and vibrational degrees of freedom (quantum harmonic oscillator), using NIST harmonic frequencies, providing physically motivated initial guesses for root-finding inversions and a temperature-dependent adiabatic gradient.
\end{enumerate}
Numerical inversions via a hybrid Newton--Brentq algorithm (for 1D inversions) or bounded least-squares with Trust Region Reflective method (for 2D $s$--$\rho$ inversions) provide the EOS in the $(s,P)$, $(\rho,T)$, $(s,\rho)$, and $(\rho,P)$ bases required by structure and evolution codes, with warm-start continuation across the composition grid for efficiency.

\begin{acknowledgments}
This work makes use of the AQUA equation of state \citep{Haldemann2020}, the CMS19 \citep{Chabrier2019} and CD21 \citep{Chabrier2021} hydrogen--helium tables, and the non-ideal H--He mixing data of \citet{Howard2023a}.
\end{acknowledgments}

\bibliographystyle{aasjournalv7}
\bibliography{references}

\end{document}
